\documentclass[11pt]{article}

% Change "review" to "final" to generate the final (sometimes called camera-ready) version.
% Change to "preprint" to generate a non-anonymous version with page numbers.
\usepackage[review]{acl}

% Standard package includes
\usepackage{times}
\usepackage{latexsym}

% For proper rendering and hyphenation of words containing Latin characters (including in bib files)
\usepackage[T1]{fontenc}
% For Vietnamese characters
% \usepackage[T5]{fontenc}
% See https://www.latex-project.org/help/documentation/encguide.pdf for other character sets

% This assumes your files are encoded as UTF8
\usepackage[utf8]{inputenc}

% This is not strictly necessary, and may be commented out,
% but it will improve the layout of the manuscript,
% and will typically save some space.
\usepackage{microtype}

% This is also not strictly necessary, and may be commented out.
% However, it will improve the aesthetics of text in
% the typewriter font.
\usepackage{inconsolata}

%Including images in your LaTeX document requires adding
%additional package(s)
\usepackage{graphicx}

\usepackage[table]{xcolor}
\usepackage{array}
\usepackage{booktabs}


% If the title and author information does not fit in the area allocated, uncomment the following
%
%\setlength\titlebox{<dim>}
%
% and set <dim> to something 5cm or larger.

% Additional packages for this draft
\usepackage{amsmath}
\usepackage{amssymb}
\usepackage{booktabs}
\usepackage{multirow}
\usepackage{xcolor}
\usepackage{url}

\newcommand{\method}{ColorRef}
\newcommand{\lab}{CIELAB}
\newcommand{\de}{\Delta E}
\newcommand{\TODO}[1]{\textcolor{red}{[TODO: #1]}}

\title{ColorRef: Interactive Reference Games for Probing the Behavioral Geometry of Color Grounding}




\author{Pablo Loyola\textsuperscript{1}, Andres Hoyos-Idobro\textsuperscript{2}, Edison Marrese-Taylor\textsuperscript{3} \\
    \textsuperscript{1} Rakuten Institute of Technology; Rakuten Group, Inc; Tokyo, Japan\\
    \textsuperscript{2} Rakuten Institute of Technology; Rakuten Group, Inc; Paris, France\\
    \textsuperscript{3} Graduate School of Engineering, The University of Tokyo\\
    \{pablo.a.loyola,andres.hoyosidrobo\}@rakuten.com, emarrese@weblab.t.u-tokyo.ac.jp
}


\begin{document}
\maketitle

\begin{abstract}
We introduce \textsc{ColorRef}, an interactive reference-game framework for
evaluating language-to-color grounding. Built from 518{,}830 English color
descriptions paired with sRGB and CIELAB targets, \textsc{ColorRef} asks
models to predict a color, receive natural-language feedback, and revise
their guess over multiple rounds. By comparing oracle and LLM teachers, the
framework separates one-shot prediction, feedback following, feedback
bandwidth, and feedback generation. Controlled oracle feedback substantially
improves accuracy and behavioral geometry: Qwen3-14B improves from 40.58 to
22.32 mean $\Delta E_{76}$, while distance-geometry alignment rises from
0.461 to 0.796. Sequential interaction outperforms same-budget single-turn
feedback, and abstract descriptions benefit most. LLM-teacher experiments
show that models can follow valid perceptual feedback better than LLMs can
generate it, exposing feedback generation as a distinct bottleneck.
\end{abstract}



\section{Introduction}
\label{sec:introduction}

Color descriptions are a compact testbed for grounded language. Some names
map directly to perceptual regions, as in \textit{dark green} or
\textit{pale blue}; others refer through objects and materials, as in
\textit{ocean blue} or \textit{mustard yellow}; and many are associative or
idiosyncratic, as in \textit{abandoned hope}, \textit{woman of the dunes},
or \textit{villainy purple}. These expressions are linguistic, but their
referents lie in a continuous perceptual space. Color therefore lets us ask
whether models can do more than associate words with labels: can they use
language to navigate perceptual structure?

Grounded communication is rarely a single prediction. People often negotiate
color reference through corrections such as \textit{lighter},
\textit{warmer}, \textit{more muted}, or \textit{less blue}. Such feedback is
linguistically simple but geometrically meaningful: it specifies local
directions in a shared perceptual space. A model that guesses poorly at first
may still be grounded if it revises in the right direction; conversely, a
model may produce fluent feedback that points another agent away from the
target. One-shot evaluation cannot distinguish these cases.

\begin{figure}[t]
    \centering
    \includegraphics[width=\columnwidth]{latex/figures/example.png}
\caption{
Paired trajectory for \textit{``very scary berry''}. Both conditions start
from the same initial guess. Oracle feedback gives geometrically valid
corrections, steering the guess toward the purple target
($\Delta E_{76}=4.4$), whereas LLM feedback is fluent but misdirected,
leaving the guess in the red region ($\Delta E_{76}=105.1$). The example
illustrates the gap between following valid feedback and generating valid
feedback.
}
    \label{fig:intro_trajectory}
\end{figure}

Most language-to-color evaluations are one-shot: given a description, a model
predicts a color, retrieves a candidate, or chooses among alternatives
\citep{mcmahan2015bayesian,monroe2017colors}. Recent work also compares
language-model representations with perceptual color geometry
\citep{abdou-etal-2021-language,loyola-etal-2023-perceptual}. These
approaches test whether initial predictions or representations align with
color space. We ask a complementary behavioral question: can models use
natural-language feedback to move through color space, and can they generate
feedback that helps another model do so?

We introduce \textsc{ColorRef}, an interactive reference-game framework for
probing color grounding in language models. In each game, a model receives a
natural-language color description and predicts an sRGB hex code. The
prediction is converted to CIELAB space and compared with a hidden target. A
teacher then provides natural-language feedback, and the model revises its
guess for up to three rounds. The resulting trajectory, illustrated in
Figure~\ref{fig:intro_trajectory}, lets us evaluate final accuracy, movement
direction, feedback-constraint satisfaction, and alignment between model
behavior and perceptual color geometry.

\begin{table}[t]
\centering
\small
\begin{tabular}{@{}lll@{}}
\toprule
Regime & Description & Swatch \\
\midrule
Explicit & faded old rose &
\cellcolor[HTML]{D2917A}\phantom{xxxx} \\
& black as coal &
\cellcolor[HTML]{1F2829}\phantom{xxxx} \\
& rich purple paint &
\cellcolor[HTML]{706CC9}\phantom{xxxx} \\
% & yellow dandelions &
% \cellcolor[HTML]{F0BA06}\phantom{xxxx} \\
\midrule
Prototype & the green of limes &
\cellcolor[HTML]{65C204}\phantom{xxxx} \\
& sweet orange sunset &
\cellcolor[HTML]{F0991A}\phantom{xxxx} \\
& montmorency cherry &
\cellcolor[HTML]{E13D2B}\phantom{xxxx} \\
% & crystal ocean shores &
% \cellcolor[HTML]{26A5C0}\phantom{xxxx} \\
\midrule
Compound & gravel grey &
\cellcolor[HTML]{8F887C}\phantom{xxxx} \\
& cryogenic electric blue &
\cellcolor[HTML]{1926FF}\phantom{xxxx} \\
& fresh kiss red &
\cellcolor[HTML]{D33750}\phantom{xxxx} \\
% & corporate blues &
% \cellcolor[HTML]{0068E3}\phantom{xxxx} \\
\midrule
 Abstract & woman of the dunes &
\cellcolor[HTML]{66615E}\phantom{xxxx} \\
& ketchup and mayo mix &
\cellcolor[HTML]{F3AC88}\phantom{xxxx} \\
% & daddy's boy blue &
%\cellcolor[HTML]{160385}\phantom{xxxx} \\
& under the seafoam blue &
\cellcolor[HTML]{47A8AB}\phantom{xxxx} \\
\bottomrule
\end{tabular}
\caption{
Sample color descriptions from the \textsc{ColorRef} corpus, grouped by
semantic regime. Swatches show crowd-aggregated ground-truth sRGB targets.
}
\label{tab:colorref_swatch_examples}
\end{table}

A key design choice is that the teacher is an experimental condition rather
than an uncontrolled conversational partner. Deterministic oracle teachers
emit controlled corrections along perceptual axes such as lightness,
red--green, yellow--blue, and saturation, isolating feedback following. We
vary bandwidth---one, two, or three constraints per turn---and compare
canonical axis-style messages with natural template paraphrases. We also
evaluate LLM teachers, which must generate feedback themselves. This
decomposes interactive grounding into initial prediction, feedback following,
bandwidth, wording, and feedback generation.

We build \textsc{ColorRef} from 518{,}830 English color descriptions paired
with sRGB colors, RGB/HSV/CIELAB coordinates, and semantic-regime labels.
Table~\ref{tab:colorref_swatch_examples} shows examples from four regimes:
\textit{explicit grounded}, \textit{prototype mediated},
\textit{compound associative}, and \textit{abstract idiosyncratic}. These
regimes test whether interaction helps differently when a description
directly names a color versus indirectly evokes one.

% We report color error as $\Delta E_{76}$, the Euclidean distance between
% predicted and target colors in CIELAB space \citep{colorimetry1986official}.
% Although values around $1$--$3$ are often associated with small or
% near-threshold differences in controlled side-by-side comparisons
% \citep{sharma2005ciede2000,khashayar2014perceptibility}, our task is more
% open-ended: models infer a crowd-sourced target from a short description.
% We therefore use $\Delta E_{76}$ as a perceptual-space distance: lower is
% better, but improved predictions may remain visually distinct from the exact
% target.

% Our experiments show that interaction substantially improves color grounding.
% On Qwen3-14B, one-shot prediction on a 12{,}000-example balanced set yields
% mean error 40.58 $\Delta E_{76}$. With three-constraint oracle feedback,
% error drops to 22.32; template feedback with the same bandwidth reaches
% 22.88. Bandwidth is the main driver in controlled ablations: moving from one
% to two constraints reduces final error from 28.13 to 22.78 for axis
% feedback, while three constraints reach 21.94. Crucially, interaction is not
% reducible to receiving more information at once: with the same total budget
% of three constraints, three one-constraint turns reach 28.13, while one
% three-constraint turn reaches only 31.88.

% Semantic-regime analysis shows that abstract descriptions are hardest in one
% shot but benefit most from interaction. Under three-constraint oracle
% feedback, abstract names improve from 49.92 to 24.05 $\Delta E_{76}$, the
% largest relative gain among regimes. Feedback also improves global geometry:
% Spearman alignment between model-output distances and true LAB distances
% rises from 0.461 in the one-shot setting to 0.796 under three-constraint
% axis feedback.

% LLM-teacher experiments reveal a second bottleneck. The guesser can follow
% well-formed feedback: a matched-vocabulary oracle teacher reaches 28.84 final
% $\Delta E$ with constraint satisfaction 0.886 on the diagnostic set. However,
% LLM-generated teachers perform substantially worse. The LAB-aware LLM teacher
% reaches 31.64, while hex-only and oracle-assisted LLM teachers end around
% 35.6 $\Delta E$. Audits and paired trajectories show that the problem is not
% mainly leakage or formatting, but directional correctness: LLM feedback is
% often fluent while pointing in the wrong perceptual direction.

We report color error as $\Delta E_{76}$, the Euclidean distance between
predicted and target colors in CIELAB space \citep{colorimetry1986official}.
Because our task asks models to infer a crowd-sourced target from a short
description, we treat $\Delta E_{76}$ as a perceptual-space distance rather
than a strict just-noticeable-difference measure: lower is better, but even
improved predictions may remain visually distinct from the exact target.

Our experiments show that interaction improves color grounding in ways that
one-shot prediction cannot capture. Controlled oracle feedback roughly halves
mean error on the 12{,}000-example scale-up set, with Qwen3-14B improving
from 40.58 to 22.32 $\Delta E_{76}$ under three-constraint feedback.
Ablations show that feedback bandwidth matters, but that interaction is not
just a larger information payload: at matched total constraint budgets,
sequential feedback outperforms delivering the same constraints in one turn.
The benefits are strongest for abstract descriptions, which are hardest in
one shot but substantially recoverable through feedback.

The same experiments reveal a teacher-side bottleneck. When feedback is
geometrically valid, the guesser can often follow it, improving both final
accuracy and global alignment with LAB geometry. LLM teachers, however,
frequently produce fluent feedback that points in the wrong perceptual
direction. Thus, \textsc{ColorRef} separates two capabilities that are often
conflated in interactive settings: using feedback to correct a trajectory,
and generating feedback that is itself geometrically correct.

Together, these results show that interactive grounding is a stronger test
than one-shot prediction: models must not only name a point in color space,
but use language to correct their trajectory through it. \textsc{ColorRef}
makes this behavior measurable, revealing a clear gap between following
perceptual feedback and generating it.

% --------------------------------------------------------------------
\section{Related Work}
\label{sec:related_work}

\paragraph{Color language, grounding, and language models.}
Color has long served as a testbed for links between language, perception,
and categorization. Classic work identifies cross-linguistic regularities in
basic color terms \citep{berlin1991basic}, while later accounts model color
naming through perceptual geometry and communicative efficiency
\citep{regier2007color,zaslavsky2018efficient}. In NLP, grounded color
semantics has been studied through probabilistic and neural models that map
linguistic expressions to perceptual observations
\citep{mcmahan2015bayesian,monroe2017colors}. Recent work asks whether
language models encode perceptual color structure without direct grounding:
\citet{abdou-etal-2021-language} compare monolexemic color-term
representations with CIELAB geometry, and
\citet{loyola-etal-2023-perceptual} show that alignment degrades with
subjectivity and abstractness. \textsc{ColorRef} extends this line from
static representations and one-shot predictions to feedback-driven
behavioral trajectories.

\paragraph{Comparatives, reference games, and pragmatics.}
Our protocol builds on comparative color language and reference games.
\citet{winn-muresan-2018-lighter} model terms such as \textit{lighter} and
\textit{darker} as directions in color space, while
\citet{monroe2017colors} study pragmatic color reference as target
identification among contextual alternatives. More broadly, reference games
and Rational Speech Act models treat communication as speaker--listener
coordination
\citep{lewis1969convention,frank2012predicting,goodman2016pragmatic}, and
grounded clarification work distinguishes perceptual from collaborative
grounding \citep{benotti2021recipe}. \textsc{ColorRef} adapts these ideas to
continuous correction: the relevant context is the model's current guess,
not a fixed distractor set.

\paragraph{Feedback, revision, and LLM teachers.}
Interactive grounding work argues that models should be evaluated not only by
task success, but by how common ground is built over time
\citep{imai2025measuring}. Feedback is also central to instruction tuning,
RLHF, AI feedback, critique, and iterative revision
\citep{ouyang2022training,bai2022constitutional,
saunders2022selfcritiquing,madaan2023selfrefine,shinn2023reflexion}. In
grounded settings, however, feedback quality is itself a source of error: a
teacher may produce fluent guidance that is not perceptually correct.
\textsc{ColorRef} therefore separates the ability to follow feedback from
the ability to generate it, using oracle and LLM teacher conditions.
% -----------------------------
% \section{The \textsc{ColorRef} Benchmark}
% \label{sec:colorref}

% \textsc{ColorRef} evaluates whether language models can ground short
% color descriptions in continuous perceptual space through interaction.
% Each example consists of a natural-language color description, a target
% sRGB color, and corresponding perceptual coordinates. A model first
% predicts a color from the description alone, then revises its prediction
% after receiving feedback. The benchmark records the full sequence of
% guesses, making it possible to evaluate both final accuracy and
% trajectory-level behavior.

% \subsection{Dataset}
% \label{sec:dataset}

% We build \textsc{ColorRef} a corpus of English
% color descriptions paired with sRGB hex codes derived from ColorNames website\footnote{\url{colornames.org}}. After preprocessing and
% filtering, the corpus contains 518{,}830 examples. Each row contains the
% description, target hex code, RGB coordinates, HSV coordinates, CIELAB
% coordinates, parser outputs, and semantic-regime annotations. We use
% CIELAB as the main evaluation space because distances in LAB are more
% perceptually meaningful than raw RGB distances.

% Each example is written as $(x_i, h_i, c_i, r_i),$, where $x_i$ is the color description, $h_i$ is the target sRGB hex code,
% $c_i \in \mathbb{R}^3$ is the corresponding LAB coordinate, and $r_i$ is
% a semantic-regime label. The guesser observes only $x_i$ at the initial
% turn; $h_i$ and $c_i$ are hidden from the guesser and used only by
% teachers and evaluators.

% \subsection{Semantic Regimes}
% \label{sec:regimes}

% Color names vary in how directly they constrain the target color. To
% analyze this variation, we assign each description to one of four
% mutually exclusive semantic regimes using a lightweight parser and
% component scores for explicitness, prototype anchoring, rarity, and
% abstraction. These labels are intended as coarse analysis categories,
% not as exhaustive semantic parses.

% \begin{table}[t]
% \centering
% \small
% \begin{tabular}{@{}lrr@{}}
% \toprule
% \textbf{Regime} & \textbf{Count} & \% \\
% \midrule
% Abstract Idiosyncratic & 238{,}416 & 46.0 \\
% Prototype Mediated & 137{,}420 & 26.5 \\
% Compound Associative & 94{,}246 & 18.2 \\
% Explicit Grounded & 48{,}748 & 9.4 \\
% \bottomrule
% \end{tabular}
% \caption{Distribution of semantic regimes in the full
% \textsc{ColorRef} corpus.}
% \label{tab:regime_distribution}
% \end{table}

% We group descriptions into four regimes: \textbf{explicit grounded}
% descriptions have direct color heads and recognized modifiers, such as
% \textit{dark green} or \textit{aged bronze}; \textbf{prototype mediated}
% descriptions are anchored by objects, materials, or natural prototypes, such
% as \textit{ocean blue} or \textit{mustard yellow}; \textbf{compound
% associative} descriptions combine a partial color anchor with additional
% associative content, as in \textit{corporate blues}; and \textbf{abstract
% idiosyncratic} descriptions are affective, metaphorical, or idiosyncratic,
% with weak direct color anchors, as in \textit{abandoned hope}.

% These regimes distinguish description validity from description
% directness. All evaluation examples are treated as plausible color
% descriptions, but they differ in how explicitly they identify a region of
% color space. This allows us to test whether interaction helps most when
% the initial linguistic signal is direct or when it is ambiguous. Details about the implementation of the parser are described in Appendix 

% \subsection{Evaluation Subsets}
% \label{sec:subsets}

% For the experiments on this paper, we construct three evaluation subsets from the full corpus. The primary
% scale-up subset, \texttt{main\_12000}, contains 12{,}000 examples, balanced
% with 3{,}000 examples per regime. The development subset,
% \texttt{main\_4000}, contains 4{,}000 examples, with 1000 examples per
% regime and additional balancing across hue, lightness, and saturation
% bins. The \texttt{debug\_4000} subset contains 4{,}000 examples, with 1000
% examples per regime, and is used only for LLM-teacher diagnostics.

% Balancing by regime prevents the evaluation from being dominated by the
% most common abstract-idiosyncratic names. Balancing by color bins reduces
% the risk that performance differences are driven by overrepresented hue,
% lightness, or saturation regions.

% \subsection{Interactive Color Reference Game}
% \label{sec:game}

% Each \textsc{ColorRef} instance is a multi-turn game between a
% \emph{guesser} and a \emph{teacher}. At turn $0$, the guesser receives only
% the description $x_i$ and predicts a hex code $\hat{h}_{i,0}=G(x_i)$, which
% we parse into sRGB and convert to LAB as
% $\hat{c}_{i,0}=\mathrm{LAB}(\hat{h}_{i,0})$. At feedback round $t$, the
% teacher observes the target $c_i$ and current guess $\hat{c}_{i,t}$, then
% emits feedback $f_{i,t}=T(c_i,\hat{c}_{i,t})$. The guesser revises its
% prediction as $\hat{h}_{i,t+1} = G(x_i, \hat{h}_{i,t}, f_{i,t})$, which is again parsed and converted to LAB. The game runs for at most
% $T=3$ feedback rounds, stopping early when
% $\lVert c_i-\hat{c}_{i,t}\rVert_2<5.0$. All main experiments use
% deterministic decoding with temperature $0$.

% \subsection{Metrics}
% \label{sec:metrics}

% The primary metric is $\Delta E_{76}$, the Euclidean distance between the
% target and predicted colors in CIELAB space: $ e_{i,t} = \lVert c_i - \hat{c}_{i,t} \rVert_2.$ Lower values indicate more perceptually similar colors, with $0$ denoting
% an exact match. We use $\Delta E_{76}$ because it is simple, reproducible,
% and directly compatible with LAB coordinates; we discuss more recent
% perceptual color-difference formulas such as CIEDE2000 as a limitation.
% We report initial error $e_{i,0}$, final error $e_{i,T}$, absolute
% improvement: $\Delta_i = e_{i,0} - e_{i,T}$, and relative improvement: $ \mathrm{RI}_i =
%     \frac{e_{i,0} - e_{i,T}}{e_{i,0} + \epsilon}$. 

% We also evaluate feedback-following behavior. \textbf{Directional
% alignment}, or directional precision measures whether the model's update moves toward the target.
% At turn $t$, there are two relevant vectors in LAB space: the ideal
% correction vector $c_i - \hat{c}_{i,t}$, pointing from the current guess
% to the target, and the model's actual update vector
% $\hat{c}_{i,t+1} - \hat{c}_{i,t}$, pointing from the current guess to the
% next guess. We compute the cosine similarity between these two vectors:

% \begin{equation}
%         A_{i,t} =
%     \frac{(c_i - \hat{c}_{i,t})^\top
%     (\hat{c}_{i,t+1} - \hat{c}_{i,t})}
%     {\lVert c_i - \hat{c}_{i,t} \rVert_2
%      \lVert \hat{c}_{i,t+1} - \hat{c}_{i,t} \rVert_2}
% \end{equation}

% This value is $1$ when the model moves exactly toward the target, $0$ when
% the update is orthogonal to the target direction, and negative when the
% model moves away from the target. Thus, directional alignment captures
% whether feedback induces a useful local movement, independently of the
% step size.



% \textbf{Constraint satisfaction} measures whether the next guess obeys the
% signed constraint expressed by feedback. For example, if the teacher says
% \textit{make it darker}, the next guess should reduce $L^\ast$; if the
% teacher says \textit{more blue}, it should move in the corresponding LAB
% direction. Unlike directional alignment, which asks whether the whole
% update points toward the target, constraint satisfaction asks whether the
% model follows the specific axis-level instruction expressed by the teacher.

% Finally, we compute \textbf{behavioral geometry alignment}. For a set of
% examples, we compare pairwise distance matrices over true LAB targets
% with pairwise distance matrices over model guesses or trajectories, using
% Spearman correlation and kernel alignment. This measures whether
% interaction makes the geometry of model behavior more consistent with
% the underlying perceptual geometry.

\section{The \textsc{ColorRef} Framework}
\label{sec:colorref}

\textsc{ColorRef} evaluates whether language models can ground short color
descriptions in continuous perceptual space through interaction. Each
example contains a natural-language color description, a target sRGB color,
and corresponding perceptual coordinates. A model first predicts a color
from the description alone, then revises its prediction after receiving
feedback. The full sequence of guesses supports evaluation of both final
accuracy and trajectory-level behavior.

\subsection{Dataset}
\label{sec:dataset}

We build \textsc{ColorRef} from English color descriptions paired with sRGB
hex codes derived from the ColorNames website\footnote{\url{colornames.org}, published  under CC0 1.0 license}.
After preprocessing and filtering, the corpus contains 518{,}830 examples.
Each row contains the description, target hex code, RGB, HSV, and CIELAB
coordinates, parser outputs, and semantic-regime annotations. We use CIELAB
as the main evaluation space because distances in LAB are more perceptually
meaningful than raw RGB distances.

Each example is written as $(x_i,h_i,c_i,r_i)$, where $x_i$ is the color
description, $h_i$ is the target sRGB hex code, $c_i \in \mathbb{R}^3$ is
the corresponding LAB coordinate, and $r_i$ is a semantic-regime label. The
guesser observes only $x_i$ at the initial turn; $h_i$ and $c_i$ are hidden
from the guesser and used only by teachers and evaluators. Details on the preprocessing steps in Appendix \ref{ap:prepro}, \ref{ap:val-filtering}.

\subsection{Semantic Regimes}
\label{sec:regimes}

Color names vary in how directly they constrain the target color. We assign
each description to one of four mutually exclusive semantic regimes using a
lightweight parser with component scores for explicitness, prototype
anchoring, rarity, and abstraction ( See details in Appendix \ref{ap:parser}). These labels are coarse analysis
categories, not exhaustive semantic parses.

\begin{table}[t]
\centering
\small
\begin{tabular}{@{}lrr@{}}
\toprule
\textbf{Regime} & \textbf{Count} & \% \\
\midrule
Abstract Idiosyncratic & 238{,}416 & 46.0 \\
Prototype Mediated & 137{,}420 & 26.5 \\
Compound Associative & 94{,}246 & 18.2 \\
Explicit Grounded & 48{,}748 & 9.4 \\
\bottomrule
\end{tabular}
\caption{Distribution of semantic regimes in the full
\textsc{ColorRef} corpus.}
\label{tab:regime_distribution}
\end{table}

% We group descriptions into four regimes: \textbf{explicit grounded}
% descriptions have direct color heads and recognized modifiers, such as
% \textit{dark green} or \textit{aged bronze}; \textbf{prototype mediated}
% descriptions are anchored by objects, materials, or natural prototypes, such
% as \textit{ocean blue} or \textit{mustard yellow}; \textbf{compound
% associative} descriptions combine a partial color anchor with additional
% associative content, as in \textit{corporate blues}; and \textbf{abstract
% idiosyncratic} descriptions are affective, metaphorical, or idiosyncratic,
% with weak direct color anchors, as in \textit{abandoned hope}. The regimes distinguish description validity from description directness:
% all evaluation examples are plausible color descriptions, but differ in how
% explicitly they identify a region of color space. This lets us test whether
% interaction helps most when the initial linguistic signal is direct or
% ambiguous. Parser details are given in Appendix~\ref{app:dataset_construction}.

We use four regimes: \textbf{explicit grounded} descriptions directly name a
color or modifier, as in \textit{dark green}; \textbf{prototype mediated}
descriptions refer through objects or materials, as in \textit{ocean blue};
\textbf{compound associative} descriptions combine a partial color anchor
with additional content, as in \textit{corporate blues}; and
\textbf{abstract idiosyncratic} descriptions are affective or metaphorical,
as in \textit{abandoned hope}. These regimes separate description validity
from directness: all examples are plausible color descriptions, but differ
in how explicitly they identify a region of color space.

\subsection{Evaluation Subsets}
\label{sec:subsets}

We construct three evaluation subsets from the full corpus.
\texttt{main\_12000} is the primary scale-up subset, with 12{,}000 examples
balanced across regimes. \texttt{main\_4000} contains 4{,}000 examples,
balanced across regimes and hue, lightness, and saturation bins.
\texttt{debug\_4000} contains 4{,}000 examples balanced across regimes and
is used only for LLM-teacher diagnostics.

Regime balancing prevents evaluation from being dominated by the most common
abstract-idiosyncratic names, while color-bin balancing reduces the risk
that differences are driven by overrepresented hue, lightness, or saturation
regions.

\subsection{Interactive Color Reference Game}
\label{sec:game}

Each \textsc{ColorRef} instance is a multi-turn game between a
\emph{guesser} and a \emph{teacher}. At turn $0$, the guesser receives only
the description $x_i$ and predicts a hex code $\hat{h}_{i,0}=G(x_i)$, which
is parsed as sRGB and converted to LAB,
$\hat{c}_{i,0}=\mathrm{LAB}(\hat{h}_{i,0})$. At feedback round $t$, the
teacher observes the target $c_i$ and current guess $\hat{c}_{i,t}$, emits
$f_{i,t}=T(c_i,\hat{c}_{i,t})$, and the guesser revises its prediction as
$\hat{h}_{i,t+1}=G(x_i,\hat{h}_{i,t},f_{i,t})$. The new prediction is again
converted to LAB. Games run for at most $T=3$ feedback rounds and stop early
when $\lVert c_i-\hat{c}_{i,t}\rVert_2<5.0$. All main experiments use
temperature-$0$ decoding.

\subsection{Metrics}
\label{sec:metrics}

The primary metric is $\Delta E_{76}$, the Euclidean distance between target
and predicted colors in CIELAB space:
$e_{i,t}=\lVert c_i-\hat{c}_{i,t}\rVert_2$. Lower values indicate more
perceptually similar colors, with $0$ denoting an exact match. We use
$\Delta E_{76}$ because it is simple, reproducible, and directly compatible
with LAB coordinates; we discuss newer perceptual color-difference formulas
such as CIEDE2000 as a limitation. We report initial error $e_{i,0}$, final
error $e_{i,T}$, absolute improvement
$\Delta_i=e_{i,0}-e_{i,T}$, and relative improvement
$\mathrm{RI}_i=(e_{i,0}-e_{i,T})/(e_{i,0}+\epsilon)$.

We also evaluate feedback-following behavior. \textbf{Directional alignment}
measures whether the model's update moves toward the target. At turn $t$, we
compare the ideal correction vector $c_i-\hat{c}_{i,t}$ with the model's
actual update vector $\hat{c}_{i,t+1}-\hat{c}_{i,t}$ using cosine similarity:
\begin{equation}
        A_{i,t} =
    \frac{(c_i - \hat{c}_{i,t})^\top
    (\hat{c}_{i,t+1} - \hat{c}_{i,t})}
    {\lVert c_i - \hat{c}_{i,t} \rVert_2
     \lVert \hat{c}_{i,t+1} - \hat{c}_{i,t} \rVert_2}.
\end{equation}
This value is $1$ when the model moves exactly toward the target, $0$ for an
orthogonal update, and negative when the model moves away.

\textbf{Constraint satisfaction} measures whether the next guess obeys the
signed constraint expressed by feedback. For example, if the teacher says
\textit{make it darker}, the next guess should reduce $L^\ast$; if the
teacher says \textit{more blue}, it should move in the corresponding LAB
direction. Unlike directional alignment, which evaluates the whole update,
constraint satisfaction checks whether the model follows the specific
axis-level instruction.

Finally, \textbf{behavioral geometry alignment} compares pairwise distance
matrices over true LAB targets with pairwise distance matrices over model
guesses or trajectories, using Spearman correlation and kernel alignment.
This measures whether interaction makes model behavior more consistent with
the underlying perceptual geometry.


%------------------------------

% \section{Experimental Setup}
% \label{sec:experimental_setup}

% We design the experiments to answer five questions: (i) whether feedback
% improves over one-shot color prediction; (ii) whether feedback bandwidth
% matters more than wording; (iii) whether sequential feedback helps beyond
% the same total amount of feedback delivered at once; (iv) whether models can
% generate useful feedback, rather than merely follow it; and (v) whether
% improvements hold across semantic regimes and model sizes.

% \subsection{Teacher Conditions}
% \label{sec:teachers}

% The teacher determines what feedback the guesser receives after each
% prediction. We consider deterministic oracle teachers and LLM teachers.
% Oracle teachers are not intended to model human tutors; they provide
% controlled feedback that isolates the guesser's ability to interpret and
% follow perceptual corrections. LLM teachers provide a more naturalistic but
% less controlled setting in which feedback must itself be generated by a
% language model.

% \paragraph{Oracle teachers.}
% Oracle teachers observe the target LAB coordinate $c_i$ and the current
% guess $\hat{c}_{i,t}$. They compute the difference along four perceptual
% axes: LAB lightness $L^\ast$, LAB red--green $a^\ast$, LAB yellow--blue
% $b^\ast$, and HSV saturation. Candidate corrections are filtered by minimum
% thresholds and ranked by normalized absolute error. The teacher then emits
% up to $k$ corrections, where $k$ is the feedback bandwidth.

% The canonical axis feedback uses short direction phrases such as
% \textit{make it lighter}, \textit{make it more blue}, or \textit{make it
% more muted}. We use three oracle variants. The \textbf{minimal oracle} emits
% exactly one canonical correction per turn. The \textbf{axis oracle} emits
% one to three canonical corrections per turn, depending on the configured
% bandwidth. The \textbf{template oracle} uses the same selected corrections as
% the axis oracle, but realizes them with deterministic natural-language
% paraphrases.

% \paragraph{LLM teachers.}
% We also evaluate LLM-generated feedback. The LLM teacher is instructed to use
% a constrained perceptual vocabulary without revealing the target hex code.
% We consider three variants: a \textbf{hex-only LLM teacher}, which sees the
% target and current guess hex codes but not explicit LAB coordinates; a
% \textbf{LAB-aware LLM teacher}, which sees both hex values and LAB
% coordinates; and an \textbf{oracle-assisted LLM teacher}, which receives
% oracle directions and paraphrases them into feedback. To separate
% teacher-generation errors from guesser-compliance errors, we include a
% matched-vocabulary oracle ceiling, \textbf{template oracle with LLM
% vocabulary}, which uses deterministic oracle directions but phrases them
% with the same constrained vocabulary available to the LLM teacher. Details of the prompts used in Appendix \ref{ap:teacher_prompts}

% \subsection{Feedback Bandwidth and Information Budget}
% \label{sec:bandwidth_budget}

% We define feedback bandwidth as the maximum number of independent
% constraints included in one teacher message: $k \in \{1,2,3\}$.
% For axis and template oracles, we run separate conditions with $k=1$,
% $k=2$, and $k=3$. This isolates whether models benefit from receiving more
% perceptual information per turn. Comparing axis and template teachers at
% the same $k$ tests whether the exact wording matters once bandwidth is
% controlled.

% We additionally run an information-budget ablation. Here, a condition is
% defined by the number of feedback rounds and the total number of constraints
% available across the whole game. This lets us compare, for example, one turn
% with three constraints against three turns with one constraint each. These
% conditions have the same total feedback budget but different temporal
% structure, testing whether iterative feedback helps beyond simply providing
% more information.

% \subsection{Models and Decoding}
% \label{sec:models}

% Our main experiments use \texttt{Qwen/Qwen3-14B} \citep{yang2025qwen3} as the guesser model. We
% also run \texttt{Qwen/Qwen3-8B} for a smaller same-family replication. All
% generations use local HuggingFace Transformers \citep{wolf2020transformers} inference with deterministic
% decoding: temperature $0$, seed $13$, and a maximum generation length of 64
% tokens. The model is instructed to output a valid hex code in the form
% \texttt{\#RRGGBB}. Outputs are parsed with a conservative hex extractor;
% parsed colors are converted to LAB before evaluation.

% The main protocol uses at most $T=3$   
% feedback rounds (Section \ref{sec:results_additional} expands on this value definition) and convergence
% threshold $\Delta E < 5.0$.  Games stop early if the current guess satisfies
% this threshold. In a separate convergence ablation, we vary the maximum
% number of feedback rounds to test whether performance plateaus after three
% turns.



% \subsection{Experimental Phases}
% \label{sec:experiment_phases}

% We organize the experiments into five phases.

% \paragraph{Phase 1: one-shot and oracle feedback.}
% We compare one-shot prediction, where the model predicts a hex code from the
% description alone, against controlled oracle feedback with up to three
% revision rounds.

% \paragraph{Phase 2: bandwidth and wording.}
% On \texttt{main\_1000}, we isolate feedback bandwidth and wording using
% seven controlled conditions: $\{\text{axis}_{c1}, \text{axis}_{c2}, \text{axis}_{c3},
%   \text{template}_{c1}, \text{template}_{c2}, \\ \text{template}_{c3},
%   \text{minimal}_{c1}\}$.


% \paragraph{Phase 3: information budget.}
% We test whether iteration matters at fixed feedback budget by comparing
% matched-budget conditions, such as one turn with three constraints versus
% three turns with one constraint each.

% \paragraph{Phase 4: LLM-teacher decomposition.}
% On \texttt{debug\_400}, we run the LLM-teacher variants and the
% matched-vocabulary oracle ceiling, then audit generated feedback by detecting
% directional constraints and checking them against the true target--guess
% geometry.

% \paragraph{Phase 5: diagnostic ablations and replication.}
% We run three diagnostic studies: a longer-turn convergence sweep varying the
% number of feedback rounds, an in-context learning calibration with
% (description, hex) demonstrations, and a model-size comparison between
% Qwen3-14B and Qwen3-8B on selected conditions.



% \subsection{Reported Quantities}
% \label{sec:reported_quantities}

% For all conditions, we report mean initial error, mean final error, relative
% improvement, convergence rate, directional alignment, and constraint
% satisfaction where applicable. For the main 12{,}000-example results and
% regime analyses, we report percentile bootstrap 95\% confidence intervals
% using 1{,}000 resamples with seed $13$.

% For the convergence ablation, we additionally report best-so-far error,
% final--best gap, and regression rate. For the information-budget ablation,
% we report total constraints issued and improvement per constraint. For ICL,
% we report results by number of demonstrations and retrieval mode. For
% LLM-teacher conditions, we report directional precision of generated
% feedback and use paired trajectory visualizations to inspect failures.

% We report results at three granularities: teacher-level aggregates,
% regime-level aggregates, and geometry-level analyses comparing pairwise
% structure in true LAB colors with pairwise structure in model guesses and
% trajectories.

\section{Experimental Setup}
\label{sec:experimental_setup}

We design experiments around five questions: whether feedback improves over
one-shot prediction; whether bandwidth matters more than wording; whether
sequential feedback helps beyond the same total information delivered at
once; whether models can generate useful feedback rather than only follow
it; and whether effects hold across semantic regimes and model sizes.

\subsection{Teacher Conditions}
\label{sec:teachers}

The teacher determines the feedback after each prediction. We compare
deterministic oracle teachers, which isolate the guesser's ability to follow
controlled perceptual corrections, with LLM teachers, which generate
feedback in a more naturalistic but less controlled setting.

\paragraph{Oracle teachers.}
Oracle teachers observe the target LAB coordinate $c_i$ and current guess
$\hat{c}_{i,t}$, compute differences along LAB lightness $L^\ast$, LAB
red--green $a^\ast$, LAB yellow--blue $b^\ast$, and HSV saturation (nominally called \emph{axis}), then
rank candidate corrections by normalized absolute error after thresholding.
The teacher emits up to $k$ corrections, where $k$ is the feedback
bandwidth.

Canonical axis feedback uses short phrases such as \textit{make it
lighter}, \textit{make it more blue}, or \textit{make it more muted}. We use
three oracle variants: a \textbf{minimal oracle} that emits exactly one
canonical correction per turn, an \textbf{axis oracle} that emits one to
three canonical corrections depending on bandwidth, and a \textbf{template
oracle} that uses the same selected corrections but realizes them with
deterministic natural-language paraphrases.

\paragraph{LLM teachers.}
We also evaluate LLM-generated feedback, instructing the teacher to use a
constrained perceptual vocabulary without revealing the target hex code. We
consider a \textbf{hex-only LLM teacher}, which sees the target and current
guess hex codes; a \textbf{LAB-aware LLM teacher}, which also sees LAB
coordinates; and an \textbf{oracle-assisted LLM teacher}, which paraphrases
oracle directions into feedback. To separate teacher-generation errors from
guesser-compliance errors, we include a matched-vocabulary oracle ceiling:
\textbf{template oracle with LLM vocabulary}, which uses deterministic oracle
directions phrased with the same vocabulary available to the LLM teacher.
Prompt details appear in Appendix~\ref{ap:teacher_prompts}.

\subsection{Feedback Bandwidth and Information Budget}
\label{sec:bandwidth_budget}

Feedback bandwidth is the maximum number of independent constraints in one
teacher message, $k \in \{1,2,3\}$. For axis and template oracles, we run
separate $k=1$, $k=2$, and $k=3$ conditions. Varying $k$ tests whether
models benefit from more perceptual information per turn; comparing axis and
template teachers at fixed $k$ tests whether wording matters once bandwidth
is controlled.

We also run an information-budget ablation, where conditions are matched by
the total number of constraints available across the game. This lets us
compare, for example, one turn with three constraints against three turns
with one constraint each, testing whether iteration helps beyond total
feedback quantity.

\subsection{Models and Decoding}
\label{sec:models}

Our main experiments use \texttt{Qwen/Qwen3-14B} \citep{yang2025qwen3} as
the guesser model, with \texttt{Qwen/Qwen3-8B} as a smaller same-family
replication. All generations use local HuggingFace Transformers
\citep{wolf2020transformers} inference with deterministic decoding:
temperature $0$, seed $13$, and maximum generation length 64. The model is
instructed to output a valid hex code in the form \texttt{\#RRGGBB}. Outputs
are parsed with a conservative hex extractor and converted to LAB before
evaluation.

The main protocol uses at most $T=3$ feedback rounds and convergence
threshold $\Delta E < 5.0$, stopping early once the threshold is reached. A
separate convergence ablation varies the maximum number of rounds to test
whether performance plateaus after three turns.

\subsection{Experimental Phases}
\label{sec:experiment_phases}

%We organize the experiments into five phases.

\paragraph{Phase 1: one-shot and oracle feedback.}
We compare one-shot prediction from the description alone against controlled
oracle feedback with up to three revision rounds.

\paragraph{Phase 2: bandwidth and wording.}
On \texttt{main\_1000}, we isolate feedback bandwidth and wording using
seven controlled oracle conditions. Here $c1$, $c2$, and $c3$ indicate the
maximum number of directional constraints per turn, while \texttt{axis} and
\texttt{template} differ only in surface wording. The conditions are
axis$_{c1}$, axis$_{c2}$, axis$_{c3}$, template$_{c1}$, template$_{c2}$,
template$_{c3}$, and minimal$_{c1}$. Examples appear in
Appendix~\ref{app:bandwidth_wording_examples}.

\paragraph{Phase 3: information budget.}
We test whether iteration matters at fixed feedback budget by comparing
matched-budget conditions, such as one turn with three constraints versus
three turns with one constraint each.

\paragraph{Phase 4: LLM-teacher decomposition.}
On \texttt{debug\_4000}, we run the LLM-teacher variants and the
matched-vocabulary oracle ceiling, then audit generated feedback by detecting
directional constraints and checking them against the true target--guess
geometry.

\paragraph{Phase 5: diagnostic ablations and replication.}
We run three diagnostic studies: a longer-turn convergence sweep, an
in-context calibration and LoRA finetuning analysis, and a model-size
comparison between Qwen3-14B and Qwen3-8B on selected conditions.

\subsection{Reported Quantities}
\label{sec:reported_quantities}

For all conditions, we report mean initial error, mean final error, relative
improvement, convergence rate, directional alignment, and constraint
satisfaction where applicable. For the main 12{,}000-example results and
regime analyses, we report percentile bootstrap 95\% confidence intervals
using 1{,}000 resamples with seed $13$. Diagnostic analyses additionally
report task-specific quantities such as best-so-far error, total constraint
budget, ICL demonstration count, and LLM-teacher directional precision.

% -----------------------------


% \section{Results}
% \label{sec:results}

% We evaluate \textsc{ColorRef} along six dimensions: improvement over one-shot prediction, feedback bandwidth and temporal structure, variation across semantic regimes, the gap between following and generating feedback, alignment with perceptual geometry, and diagnostic ablations.

% \subsection{Controlled feedback improves color grounding}
% \label{sec:results_main}

% We first compare one-shot prediction against controlled feedback on the 12{,}000-example scale-up set. In the one-shot baseline, the model receives only the color description and predicts a hex code. In feedback conditions, the model may revise its prediction over up to three feedback rounds.

% \begin{table}[t]
% \centering
% \small
% \scalebox{0.93}{
% \begin{tabular}{@{}lcccc@{}}
% \toprule
% \bf Condition & $N$ & \bf Init $\Delta E$ & \bf Final $\Delta E$ & \bf Rel. imp. \\
% \midrule
% one-shot & 12,000 & 40.58 & 40.58 & 0.000 \\
% axis$_{c3}$ & 12,000 & 40.58 & \textbf{22.32} & \textbf{0.262} \\
% template$_{c3}$ & 12,000 & 40.58 & 22.88 & 0.243 \\
% minimal$_{c1}$ & 12,000 & 40.58 & 28.29 & 0.124 \\
% template-vocab & 12,000 & 40.58 & 29.68 & 0.120 \\
% \bottomrule
% \end{tabular}
% }
% \caption{Main 12k results for Qwen3-14B. Controlled feedback substantially reduces LAB error relative to one-shot prediction.}
% \label{tab:main_results}
% \end{table}

% Table~\ref{tab:main_results} shows that controlled feedback substantially improves color grounding. One-shot prediction yields mean LAB error 40.58. With three-constraint axis feedback, final error drops to 22.32, a relative improvement of 26.2\%. Template feedback with the same bandwidth reaches a similar final error of 22.88. Single-constraint feedback also improves over one-shot, but is substantially weaker, ending around 28--30 $\Delta E$.

% This effect is stable at scale. The 95\% bootstrap interval for axis$_{c3}$ final error is [21.85, 22.76], and the corresponding interval for template$_{c3}$ is [22.38, 23.36]. Thus, interaction does not merely improve a small development set; it produces a robust reduction in perceptual error on the larger balanced evaluation set.

% Qualitatively, the trajectories show that feedback can correct large semantic misfires. For example, an abstract description may initially induce a color from the wrong region of space, but successive feedback such as \textit{more muted}, \textit{more yellow}, or \textit{more green} can move the model toward the target. This illustrates the central phenomenon: even when the initial association is wrong, local perceptual feedback can steer the model toward the intended color.

% \subsection{Bandwidth matters, but iteration matters too}
% \label{sec:results_bandwidth_budget}

% We next isolate two properties of feedback: its \emph{bandwidth}, meaning the number of independent constraints per turn, and its \emph{temporal structure}, meaning whether constraints are delivered all at once or sequentially across turns.

% \begin{table}[t]
% \centering
% \small
% \scalebox{0.88}{
% \begin{tabular}{@{}lccccc@{}}
% \toprule
% \bf Condition & \bf Max $c$ & \bf Final $\Delta E$ & \bf Rel. imp. & \bf C-sat & \bf Align. \\
% \midrule
% axis$_{c1}$ & 1 & 28.13 & 0.127 & 0.935 & 0.549 \\
% axis$_{c2}$ & 2 & 22.78 & 0.266 & 0.936 & 0.727 \\
% axis$_{c3}$ & 3 & \textbf{21.94} & 0.254 & 0.898 & 0.742 \\
% template$_{c1}$ & 1 & 29.97 & 0.128 & 0.889 & 0.483 \\
% template$_{c2}$ & 2 & 23.40 & 0.255 & 0.922 & 0.703 \\
% template$_{c3}$ & 3 & 22.32 & \textbf{0.266} & 0.876 & 0.712 \\
% minimal$_{c1}$ & 1 & 28.13 & 0.127 & 0.935 & 0.549 \\
% \bottomrule
% \end{tabular}
% }
% \caption{Bandwidth ablation on \texttt{main\_4000}. Increasing feedback from one to two constraints gives the largest gain; axis and template wording are similar at matched bandwidth.}
% \label{tab:bandwidth}
% \end{table}

% Table~\ref{tab:bandwidth} shows that bandwidth is the dominant factor among oracle-feedback variants. Moving from one to two constraints reduces final error by roughly 5--7 $\Delta E$: axis feedback improves from 28.13 to 22.78, while template feedback improves from 29.97 to 23.40. Moving from two to three constraints gives only a small additional gain. This suggests that much of the useful corrective signal can be captured by two well-chosen perceptual directions per turn.

% By contrast, wording has a smaller effect once bandwidth is controlled. At $c=3$, axis and template feedback are nearly tied: 21.94 versus 22.32 final $\Delta E$. At $c=2$, they are also close: 22.78 versus 23.40. The larger difference appears at $c=1$, where template feedback trails axis feedback by about 1.8 $\Delta E$. Overall, the model is not merely relying on fixed phrases such as \textit{make it darker}; it can use paraphrased feedback, but it benefits most from receiving enough independent constraints.

% \begin{table}[t]
% \centering
% \small
% \begin{tabular}{@{}lcccc@{}}
% \toprule
% \bf Condition & \bf Rounds & \bf $c$/turn & \bf Budget & \bf Final $\Delta E$ \\
% \midrule
% one-turn$_{c2}$ & 1 & 2 & 2 & 31.98 \\
% two-turn$_{c1}$ & 2 & 1 & 2 & \textbf{29.71} \\
% \midrule
% one-turn$_{c3}$ & 1 & 3 & 3 & 31.88 \\
% three-turn$_{c1}$ & 3 & 1 & 3 & \textbf{28.13} \\
% \midrule
% three-turn$_{c3}$ & 3 & 3 & 9 & \textbf{21.94} \\
% \bottomrule
% \end{tabular}
% \caption{Information-budget ablation on \texttt{main\_4000}. At matched total constraint budgets, iterative feedback outperforms delivering the same number of constraints in a single turn.}
% \label{tab:information_budget}
% \end{table}

% The bandwidth result raises a further question: does feedback help simply because it gives the model more information, or because the information is delivered interactively? Table~\ref{tab:information_budget} shows that the temporal structure matters. With the same total budget of three constraints, receiving one constraint per turn over three turns reaches final error 28.13, while receiving all three constraints in a single turn reaches only 31.88. The same pattern holds for budget two: two one-constraint turns reach 29.71, while one two-constraint turn reaches 31.98.

% Thus, feedback is not merely an information payload. Iterative feedback is more effective because later corrections are conditioned on the model's revised guess. The strongest condition still combines both factors: three turns with three constraints per turn reaches 21.94. This indicates that total information and sequential adaptation are complementary.

% \subsection{Semantic regimes reveal difficulty and recoverability}
% \label{sec:results_regime}

% The semantic-regime taxonomy lets us ask whether different kinds of color descriptions benefit differently from interaction. Table~\ref{tab:regime_axis} reports results for axis$_{c3}$ on the 12{,}000-example scale-up set.

% \begin{table}[t]
% \centering
% \small
% \begin{tabular}{@{}lcccc@{}}
% \toprule
% \bf Regime & \bf $N$ & \bf Init $\Delta E$ & \bf Final $\Delta E$ & \bf Rel. imp. \\
% \midrule
% abstract & 3,000 & 49.92 & 24.05 & \textbf{0.343} \\
% compound & 3,000 & 36.87 & 21.96 & 0.237 \\
% explicit & 3,000 & 36.02 & \textbf{21.15} & 0.246 \\
% prototype & 3,000 & 39.51 & 22.14 & 0.222 \\
% \bottomrule
% \end{tabular}
% \caption{Regime effects for axis$_{c3}$ on \texttt{main\_12000}. Abstract descriptions have the highest initial error but the largest relative improvement.}
% \label{tab:regime_axis}
% \end{table}

% The strongest pattern is that abstract-idiosyncratic descriptions are hardest initially but most recoverable under feedback. Their one-shot error is 49.92, far above the other regimes, but final error drops to 24.05, yielding the largest relative improvement, 34.3\%. Explicit-grounded names start much closer to the target, with initial error 36.02, and end with the lowest final error, 21.15. Prototype-mediated and compound-associative descriptions fall between these cases.

% This result complicates a simple ``abstract names are noise'' interpretation. Abstract descriptions are weakly grounded in the one-shot setting, but they are not uniformly unrecoverable. When the model receives informative perceptual feedback, many abstract examples move substantially toward the target. Interaction therefore exposes a distinction between initial ambiguity and interactive recoverability.

% The same pattern holds across teacher conditions. In the full regime-by-teacher table, abstract descriptions remain the hardest under one-shot prediction, but axis$_{c3}$ and template$_{c3}$ reduce their error by roughly half. Single-constraint teachers help less, especially for abstract names, suggesting that descriptions with broader initial uncertainty require higher-bandwidth feedback to localize the intended region.

% \subsection{Following feedback is easier than generating feedback}
% \label{sec:results_llm_teacher}

% The oracle experiments isolate the guesser's ability to follow correct feedback. We next ask whether LLMs can generate such feedback themselves. Table~\ref{tab:llm_teachers} compares LLM-teacher variants on the 4000-example diagnostic subset.

% \begin{table}[t]
% \centering
% \small
% \begin{tabular}{@{}l@{\hspace{0.2cm}} c@{\hspace{0.2cm}} c@{\hspace{0.2cm}} c@{\hspace{0.2cm}} c@{\hspace{0.2cm}}@{}}
% \toprule
% \bf Teacher & \bf $\Delta E$ & \bf Rel. imp. & \bf C-sat & \bf DP \\
% \midrule
% template-vocab oracle & \textbf{28.84} & \textbf{0.118} & \textbf{0.886} & -- \\
% LLM lab-aware & 31.64 & 0.091 & 0.456 & 0.39 \\
% LLM hex-only & 35.56 & -0.049 & 0.487 & 0.31 \\
% LLM oracle-assisted & 35.66 & -0.083 & 0.651 & 0.11 \\
% \bottomrule
% \end{tabular}
% \caption{LLM-teacher decomposition on \texttt{debug\_400}. The matched-vocabulary oracle ceiling shows that the guesser can follow well-formed feedback; LLM teachers often fail to generate geometrically correct feedback. In the table, $\Delta E$ is short for Final $\Delta E$, and DP is short directional precision.}
% \label{tab:llm_teachers}
% \end{table}

% The matched-vocabulary oracle ceiling performs best among these diagnostic conditions, with final error 28.84 and constraint satisfaction 0.886. This condition uses deterministic oracle directions but expresses them using the vocabulary available to LLM teachers. Its success shows that the guesser can use this style of feedback when the underlying direction is correct.

% The LLM teachers are substantially weaker. The LAB-aware LLM teacher performs best among generated-feedback variants, reaching final error 31.64 and positive relative improvement. The hex-only teacher performs worse, with negative relative improvement, suggesting that hex values alone are not sufficient for reliable feedback generation. Surprisingly, the oracle-assisted LLM teacher also performs poorly. Although this variant has higher constraint satisfaction than the other LLM teachers, its audited directional precision is only 0.11, indicating that the generated text often expresses the wrong geometric direction.

% % \begin{figure}[t]
% %     \centering
% %     \includegraphics[width=\columnwidth]{figures/oracle_vs_llm_teacher_pair_grid.pdf}
% %     \caption{
% %     \textbf{Oracle and LLM teachers can induce opposite trajectories.}
% %     Paired examples compare template-oracle feedback with LLM-generated feedback for the same or nearly same initial guess. Oracle feedback moves the model toward the target, while LLM feedback is often fluent but directionally incorrect, causing the trajectory to diverge in CIELAB space.
% %     }
% %     \label{fig:oracle_llm_pairs}
% % \end{figure}

% % Figure~\ref{fig:oracle_llm_pairs} illustrates the failure mode. Starting from the same or nearly same initial state, oracle feedback can move the guesser toward the target while LLM-generated feedback sends it away. The selected paired cases show that the gap is semantic and geometric, not merely a formatting issue: the teacher's language sounds plausible, but its implied direction in color space is wrong.

% These results support a decomposition of interactive grounding into two capabilities. The guesser can often follow clear perceptual feedback, as shown by oracle teachers. But generating feedback requires translating a target--guess difference into geometrically valid natural language, and this remains difficult for LLM teachers. Failure is not primarily due to leakage or output formatting; the main failure mode is fluent but misdirected feedback.

% \subsection{Feedback improves behavioral geometry alignment}
% \label{sec:results_geometry}

% The previous results evaluate pointwise error. We now ask whether feedback also improves the global geometry of model behavior. For each condition, we compute pairwise distance matrices over true LAB targets and over model guesses or trajectories, then compare them using Spearman correlation and kernel alignment.

% \begin{table}[t]
% \centering
% \small
% \scalebox{0.95}{
% \begin{tabular}{@{}lcccc@{}}
% \toprule
% \bf Condition & \bf $\rho$ init & \bf $\rho$ final & \bf  $\rho$ traj. & \bf Kernel final \\
% \midrule
% one-shot & 0.461 & 0.461 & 0.461 & 0.508 \\
% axis$_{c3}$ & 0.461 & \textbf{0.796} & \textbf{0.804} & \textbf{0.832} \\
% template$_{c3}$ & 0.461 & 0.764 & 0.789 & 0.805 \\
% minimal$_{c1}$ & 0.461 & 0.693 & 0.696 & 0.717 \\
% template-vocab & 0.461 & 0.663 & 0.666 & 0.681 \\
% \bottomrule
% \end{tabular}
% }
% \caption{Behavioral geometry alignment on \texttt{main\_12000}. Feedback makes the geometry of model guesses and trajectories more aligned with true LAB geometry.}
% \label{tab:geometry}
% \end{table}

% Table~\ref{tab:geometry} shows that one-shot guesses have moderate alignment with true LAB geometry, with Spearman $\rho=0.461$. Feedback substantially increases this alignment. Axis$_{c3}$ reaches final-guess alignment $\rho=0.796$ and trajectory-state alignment $\rho=0.804$; template$_{c3}$ follows closely with final alignment $\rho=0.764$ and trajectory alignment $\rho=0.789$. Single-constraint feedback also improves alignment, but less strongly.

% This result supports the behavioral-geometry framing of \textsc{ColorRef}. Feedback does not merely reduce the average distance between individual guesses and targets. It reorganizes the set of model outputs so that pairwise relationships among guesses better match pairwise relationships among true colors. In other words, interaction makes the model's behavior more geometrically faithful to the underlying perceptual space.

% \subsection{Additional ablations: convergence, ICL, finetuning and model size}
% \label{sec:results_additional}

% We include three additional ablations to clarify the scope of the main results.

% \paragraph{Longer games.}
% The main protocol uses three feedback rounds. A longer-turn sweep shows that this is a strong compute--quality tradeoff. Final error improves rapidly from one to three turns, from 31.88 to 21.94, but then mostly plateaus: five turns reaches 21.43 and ten turns reaches 21.04. Best-so-far error continues to improve, from 15.67 at three turns to 10.17 at ten turns, but the final--best gap grows from 6.27 to 10.87. Thus, longer games discover better intermediate colors but increasingly overshoot them. Details of the experiments in Appendix \ref{ap:longer-turn}.

% \paragraph{In-context calibration and finetuning.}
% We also test whether examples of (description, hex) pairs improve performance. ICL modestly improves one-shot prediction, reducing mean error from 41.47 to roughly 38.5 for $k=2$ examples. The gain is largest for abstract-idiosyncratic names. However, ICL provides little additional benefit once oracle feedback is available: axis$_{c3}$ without ICL reaches 21.94, while ICL feedback variants remain around 21.8--22.3. This suggests that demonstrations calibrate the initial prior, whereas per-example feedback provides the stronger localization signal Details of this experiments are shown in Appendix \ref{ap:icl-experiments}. We also tested lightweight finetuning using LoRA\cite{hu2022lora} on a balanced 4k subset of
% the data. The resulting differences were not substantial: performance
% remained close to the corresponding non-finetuned and ICL conditions, with
% no clear change in the overall conclusions. We therefore treat LoRA
% finetuning as a diagnostic intervention: small-scale parameter-efficient
% supervision provides limited additional calibration, whereas per-example
% perceptual feedback remains the stronger localization signal.

% \paragraph{Model size.}
% Finally, we compare Qwen3-14B and Qwen3-8B on one-shot, axis$_{c3}$, and template$_{c1}$ conditions. The smaller model follows the same qualitative ordering, but benefits less from feedback: Qwen3-8B improves from 44.18 to 27.99 under axis$_{c3}$, while Qwen3-14B improves from 41.47 to 21.94. Because this comparison uses only two models from the same family, we treat it as a replication of the qualitative pattern rather than a broad scaling claim. Details of this experiments are shown in Appendix \ref{ap:model-size}.

% \subsection{Summary}
% \label{sec:results_summary}

% Across these analyses, four findings stand out. First, interactive feedback substantially improves color grounding and behavioral geometry alignment. Second, bandwidth matters, but interaction is not reducible to total information: at matched constraint budgets, sequential feedback outperforms single-turn feedback. Third, abstract descriptions are difficult in one shot but interactively recoverable. Fourth, following perceptual feedback and generating perceptually valid feedback are distinct capabilities: the guesser can comply with well-formed corrections, but LLM teachers often generate fluent feedback that points in the wrong direction.


\section{Results}
\label{sec:results}

We evaluate \textsc{ColorRef} along six dimensions: improvement over one-shot
prediction, feedback bandwidth and temporal structure, semantic-regime
variation, LLM-teacher feedback generation, behavioral geometry, and
diagnostic ablations.

\subsection{Controlled feedback improves color grounding}
\label{sec:results_main}

We first compare one-shot prediction against controlled feedback on the
12{,}000-example scale-up set. The one-shot baseline predicts a hex code
from the description alone; feedback conditions allow up to three revision
rounds.

\begin{table}[t]
\centering
\small
\scalebox{0.93}{
\begin{tabular}{@{}lcccc@{}}
\toprule
\bf Condition & $N$ & \bf Init $\Delta E$ & \bf Final $\Delta E$ & \bf Rel. imp. \\
\midrule
one-shot & 12,000 & 40.58 & 40.58 & 0.000 \\
axis$_{c3}$ & 12,000 & 40.58 & \textbf{22.32} & \textbf{0.262} \\
template$_{c3}$ & 12,000 & 40.58 & 22.88 & 0.243 \\
minimal$_{c1}$ & 12,000 & 40.58 & 28.29 & 0.124 \\
template-vocab & 12,000 & 40.58 & 29.68 & 0.120 \\
\bottomrule
\end{tabular}
}
\caption{Main 12k results for Qwen3-14B. Controlled feedback substantially
reduces LAB error relative to one-shot prediction.}
\label{tab:main_results}
\end{table}

Table~\ref{tab:main_results} shows that controlled feedback substantially
improves color grounding. One-shot prediction yields mean LAB error 40.58.
With three-constraint axis feedback, final error drops to 22.32, a relative
improvement of 26.2\%; template feedback with the same bandwidth reaches
22.88. Single-constraint feedback also improves over one-shot but is weaker,
ending around 28--30 $\Delta E$.

The effect is stable at scale: the 95\% bootstrap interval for axis$_{c3}$
final error is [21.85, 22.76], and for template$_{c3}$ it is [22.38, 23.36].
Qualitatively, trajectories show that even when the initial association is
wrong, local perceptual feedback can steer the model toward the intended
color.

\subsection{Bandwidth matters, but iteration matters too}
\label{sec:results_bandwidth_budget}

We next isolate feedback \emph{bandwidth}, the number of independent
constraints per turn, and \emph{temporal structure}, whether constraints are
delivered at once or sequentially.

\begin{table}[t]
\centering
\small
\scalebox{0.88}{
\begin{tabular}{@{}lccccc@{}}
\toprule
\bf Condition & \bf Max $c$ & \bf Final $\Delta E$ & \bf Rel. imp. & \bf C-sat & \bf Align. \\
\midrule
axis$_{c1}$ & 1 & 28.13 & 0.127 & 0.935 & 0.549 \\
axis$_{c2}$ & 2 & 22.78 & 0.266 & 0.936 & 0.727 \\
axis$_{c3}$ & 3 & \textbf{21.94} & 0.254 & 0.898 & 0.742 \\
template$_{c1}$ & 1 & 29.97 & 0.128 & 0.889 & 0.483 \\
template$_{c2}$ & 2 & 23.40 & 0.255 & 0.922 & 0.703 \\
template$_{c3}$ & 3 & 22.32 & \textbf{0.266} & 0.876 & 0.712 \\
minimal$_{c1}$ & 1 & 28.13 & 0.127 & 0.935 & 0.549 \\
\bottomrule
\end{tabular}
}
\caption{Bandwidth ablation on \texttt{main\_4000}. Increasing feedback
from one to two constraints gives the largest gain; axis and template
wording are similar at matched bandwidth.}
\label{tab:bandwidth}
\end{table}

Table~\ref{tab:bandwidth} shows that bandwidth is the dominant factor among
oracle-feedback variants. Moving from one to two constraints reduces final
error by roughly 5--7 $\Delta E$: axis feedback improves from 28.13 to
22.78, while template feedback improves from 29.97 to 23.40. Moving from two
to three constraints yields a smaller additional gain. Wording has a weaker
effect once bandwidth is controlled: at $c=3$, axis and template feedback
are close, 21.94 versus 22.32 final $\Delta E$.

\begin{table}[t]
\centering
\small
\begin{tabular}{@{}lcccc@{}}
\toprule
\bf Condition & \bf Rounds & \bf $c$/turn & \bf Budget & \bf Final $\Delta E$ \\
\midrule
one-turn$_{c2}$ & 1 & 2 & 2 & 31.98 \\
two-turn$_{c1}$ & 2 & 1 & 2 & \textbf{29.71} \\
\midrule
one-turn$_{c3}$ & 1 & 3 & 3 & 31.88 \\
three-turn$_{c1}$ & 3 & 1 & 3 & \textbf{28.13} \\
\midrule
three-turn$_{c3}$ & 3 & 3 & 9 & \textbf{21.94} \\
\bottomrule
\end{tabular}
\caption{Information-budget ablation on \texttt{main\_4000}. At matched
total constraint budgets, iterative feedback outperforms delivering the same
number of constraints in a single turn.}
\label{tab:information_budget}
\end{table}

Table~\ref{tab:information_budget} shows that temporal structure also
matters. With the same total budget of three constraints, three
one-constraint turns reach 28.13 final error, while one three-constraint turn
reaches only 31.88. The same pattern holds at budget two: two
one-constraint turns reach 29.71, compared with 31.98 for one
two-constraint turn. Thus, feedback is not merely an information payload:
later corrections are useful because they are conditioned on the model's
revised guess.

\subsection{Semantic regimes reveal difficulty and recoverability}
\label{sec:results_regime}

The semantic-regime taxonomy lets us ask whether different kinds of color
descriptions benefit differently from interaction. Table~\ref{tab:regime_axis}
reports axis$_{c3}$ results on \texttt{main\_12000}.

\begin{table}[t]
\centering
\small
\begin{tabular}{@{}lcccc@{}}
\toprule
\bf Regime & \bf $N$ & \bf Init $\Delta E$ & \bf Final $\Delta E$ & \bf Rel. imp. \\
\midrule
abstract & 3,000 & 49.92 & 24.05 & \textbf{0.343} \\
compound & 3,000 & 36.87 & 21.96 & 0.237 \\
explicit & 3,000 & 36.02 & \textbf{21.15} & 0.246 \\
prototype & 3,000 & 39.51 & 22.14 & 0.222 \\
\bottomrule
\end{tabular}
\caption{Regime effects for axis$_{c3}$ on \texttt{main\_12000}. Abstract
descriptions have the highest initial error but the largest relative
improvement.}
\label{tab:regime_axis}
\end{table}

Abstract-idiosyncratic descriptions are hardest initially but most
recoverable under feedback. Their one-shot error is 49.92, far above the
other regimes, but final error drops to 24.05, yielding the largest relative
improvement, 34.3\%. Explicit-grounded names start closer to the target and
end with the lowest final error, 21.15, while prototype-mediated and
compound-associative descriptions fall between these cases.

This complicates a simple ``abstract names are noise'' interpretation.
Abstract descriptions are weakly grounded in one shot, but many become
recoverable with informative perceptual feedback. The full regime-by-teacher
table shows the same pattern: axis$_{c3}$ and template$_{c3}$ reduce
abstract-name error by roughly half, whereas single-constraint teachers help
less.

\subsection{Following feedback is easier than generating feedback}
\label{sec:results_llm_teacher}

Oracle experiments isolate the guesser's ability to follow correct feedback.
We next ask whether LLMs can generate such feedback themselves.
Table~\ref{tab:llm_teachers} compares LLM-teacher variants on
\texttt{debug\_4000}.

\begin{table}[t]
\centering
\small
\begin{tabular}{@{}l@{\hspace{0.2cm}} c@{\hspace{0.2cm}} c@{\hspace{0.2cm}} c@{\hspace{0.2cm}} c@{\hspace{0.2cm}}@{}}
\toprule
\bf Teacher & \bf $\Delta E$ & \bf Rel. imp. & \bf C-sat & \bf DP \\
\midrule
template-vocab oracle & \textbf{28.84} & \textbf{0.118} & \textbf{0.886} & -- \\
LLM lab-aware & 31.64 & 0.091 & 0.456 & 0.39 \\
LLM hex-only & 35.56 & -0.049 & 0.487 & 0.31 \\
LLM oracle-assisted & 35.66 & -0.083 & 0.651 & 0.11 \\
\bottomrule
\end{tabular}
\caption{LLM-teacher decomposition on \texttt{debug\_4000}. The
matched-vocabulary oracle ceiling shows that the guesser can follow
well-formed feedback; LLM teachers often fail to generate geometrically
correct feedback. Here, $\Delta E$ denotes final $\Delta E$, and DP denotes
directional precision.}
\label{tab:llm_teachers}
\end{table}

The matched-vocabulary oracle ceiling performs best, with final error 28.84
and constraint satisfaction 0.886. Because it uses deterministic oracle
directions phrased with the vocabulary available to LLM teachers, its success
shows that the guesser can use this feedback style when the underlying
direction is correct.

LLM teachers are substantially weaker. The LAB-aware teacher performs best
among generated-feedback variants, reaching final error 31.64 and positive
relative improvement. The hex-only teacher performs worse, suggesting that
hex values alone are insufficient for reliable feedback generation. The
oracle-assisted LLM teacher also performs poorly: despite higher constraint
satisfaction, its audited directional precision is only 0.11, indicating
that paraphrased feedback often expresses the wrong geometric direction.

These results decompose interactive grounding into two capabilities. The
guesser can often follow clear perceptual feedback, but generating feedback
requires translating a target--guess difference into geometrically valid
language. The main LLM-teacher failure mode is therefore not leakage or
formatting, but fluent, misdirected feedback.

\subsection{Feedback improves behavioral geometry alignment}
\label{sec:results_geometry}

Beyond pointwise error, we ask whether feedback improves the global geometry
of model behavior. For each condition, we compare pairwise distance matrices
over true LAB targets with pairwise distance matrices over model guesses or
trajectories, using Spearman correlation and kernel alignment.

\begin{table}[t]
\centering
\small
\scalebox{0.95}{
\begin{tabular}{@{}lcccc@{}}
\toprule
\bf Condition & \bf $\rho$ init & \bf $\rho$ final & \bf  $\rho$ traj. & \bf Kernel final \\
\midrule
one-shot & 0.461 & 0.461 & 0.461 & 0.508 \\
axis$_{c3}$ & 0.461 & \textbf{0.796} & \textbf{0.804} & \textbf{0.832} \\
template$_{c3}$ & 0.461 & 0.764 & 0.789 & 0.805 \\
minimal$_{c1}$ & 0.461 & 0.693 & 0.696 & 0.717 \\
template-vocab & 0.461 & 0.663 & 0.666 & 0.681 \\
\bottomrule
\end{tabular}
}
\caption{Behavioral geometry alignment on \texttt{main\_12000}. Feedback
makes the geometry of model guesses and trajectories more aligned with true
LAB geometry.}
\label{tab:geometry}
\end{table}

Table~\ref{tab:geometry} shows that one-shot guesses have moderate alignment
with true LAB geometry, with Spearman $\rho=0.461$. Feedback substantially
increases this alignment: axis$_{c3}$ reaches final-guess alignment
$\rho=0.796$ and trajectory-state alignment $\rho=0.804$, while
template$_{c3}$ reaches $\rho=0.764$ and $\rho=0.789$. Single-constraint
feedback also improves alignment, but less strongly. Thus, feedback not only
reduces individual errors; it makes pairwise relationships among model
outputs more faithful to perceptual color space.

% \subsection{Additional ablations: convergence, ICL, finetuning, and model size}
% \label{sec:results_additional}

% We include three additional ablations to clarify the scope of the main
% results.

% \paragraph{Longer games.}
% The main protocol uses three feedback rounds. A longer-turn sweep shows that
% final error improves rapidly from one to three turns, from 31.88 to 21.94,
% but then mostly plateaus: five turns reaches 21.43 and ten turns reaches
% 21.04. Best-so-far error continues to improve, from 15.67 at three turns to
% 10.17 at ten turns, but the final--best gap grows from 6.27 to 10.87. Thus,
% longer games discover better intermediate colors but increasingly overshoot
% them. Details appear in Appendix~\ref{ap:longer-turn}.

% \paragraph{In-context calibration and finetuning.}
% We also test whether examples of (description, hex) pairs improve
% performance. ICL modestly improves one-shot prediction, reducing mean error
% from 41.47 to roughly 38.5 for $k=2$ examples, with the largest gain for
% abstract-idiosyncratic names. Once oracle feedback is available, however,
% ICL adds little: axis$_{c3}$ without ICL reaches 21.94, while ICL feedback
% variants remain around 21.8--22.3. This suggests that demonstrations
% calibrate the initial prior, whereas per-example feedback provides the
% stronger localization signal. Details appear in
% Appendix~\ref{ap:icl-experiments}.

% We also tested parameter-efficient finetuning with LoRA~\citep{hu2021lora}
% on a balanced 4k subset of the data. The differences were not substantial:
% performance remained close to the corresponding non-finetuned and ICL
% conditions, with no clear change in the overall conclusions. We therefore
% treat LoRA finetuning as a diagnostic intervention: small-scale supervision
% provides limited additional calibration, while per-example perceptual
% feedback remains the stronger localization signal.

% \paragraph{Model size.}
% Finally, we compare Qwen3-14B and Qwen3-8B on one-shot, axis$_{c3}$, and
% template$_{c1}$ conditions. The smaller model follows the same qualitative
% ordering but benefits less from feedback: Qwen3-8B improves from 44.18 to
% 27.99 under axis$_{c3}$, while Qwen3-14B improves from 41.47 to 21.94.
% Because this comparison uses only two same-family models, we treat it as a
% replication of the qualitative pattern rather than a broad scaling claim.
% Details appear in Appendix~\ref{ap:model-size}.


\subsection{Additional ablations: convergence, ICL, finetuning, and model size}
\label{sec:results_additional}

Additional ablations clarify the scope of the main results. First, a
longer-turn sweep shows that final error improves rapidly from one to three
turns and then mostly plateaus: five and ten turns reach 21.43 and 21.04
final $\Delta E$, while the final--best gap grows, indicating overshooting
(Appendix~\ref{ap:longer-turn}). Second, in-context examples modestly
improve one-shot prediction, especially for abstract-idiosyncratic names,
but add little once oracle feedback is available: feedback variants remain
around 21.8--22.3 final $\Delta E$ (Appendix~\ref{ap:icl-experiments}).
Parameter-efficient LoRA finetuning on a balanced 4k subset likewise yields
no substantial change, suggesting that per-example perceptual feedback is a
stronger localization signal than small-scale calibration. Finally, Qwen3-8B
shows the same qualitative ordering as Qwen3-14B but benefits less from
feedback, so we treat the model-size comparison as a same-family replication
rather than a scaling claim (Appendix~\ref{ap:model-size}).

% \subsection{Summary}
% \label{sec:results_summary}

% Across analyses, four findings stand out. First, interactive feedback
% improves both color accuracy and behavioral geometry. Second, bandwidth
% matters, but sequential feedback helps beyond total information. Third,
% abstract descriptions are difficult in one shot but interactively
% recoverable. Fourth, following valid perceptual feedback and generating it
% are distinct capabilities: LLM teachers often produce fluent feedback that
% points in the wrong direction.

% Together, these results show that interactive feedback improves both pointwise
% accuracy and global behavior, while exposing a sharp gap between following
% valid perceptual feedback and generating it.


% ------------------

% \section{Discussion}
% \label{sec:discussion}

% \paragraph{Interaction as adaptive grounding.}
% The information-budget ablation shows that feedback is not merely an
% additional information channel. At matched total constraint budgets,
% sequential feedback outperforms delivering the same constraints in a single
% turn, suggesting that interaction helps because each teacher message is
% conditioned on the model's previous guess. In this sense, \textsc{ColorRef}
% tests whether language can act as a control signal for navigating perceptual
% space. A correction such as \textit{more blue} is useful only relative to the
% current guess; once the guess changes, the next useful correction changes as
% well.

% \paragraph{Why three turns?}
% The longer-turn sweep supports three feedback rounds as the main protocol.
% Final error improves rapidly from one to three turns and then largely
% plateaus. Longer games still reduce best-so-far error, but the growing gap
% between final and best-so-far error indicates overshooting. Three turns
% therefore provide a practical compute--quality tradeoff, while longer games
% are better suited for studying stopping policies or best-state selection.

% \paragraph{Demonstrations calibrate priors, feedback localizes targets.}
% In-context examples modestly improve one-shot predictions, especially for
% abstract-idiosyncratic descriptions, but add little once oracle feedback is
% available. This suggests that demonstrations calibrate the initial color
% prior, whereas per-example feedback supplies the stronger localization
% signal. It also indicates that the main \textsc{ColorRef} effects are not
% simply due to unfamiliar output format or dataset style.

% \paragraph{Following and generating feedback are distinct capabilities.}
% The LLM-teacher results reveal a sharp asymmetry: the guesser can often
% follow well-formed perceptual corrections, but LLM teachers frequently
% generate fluent feedback that points in the wrong geometric direction. This
% is a controlled form of common-ground misalignment: the teacher produces a
% plausible utterance that is not aligned with the perceptual relation needed
% for update \citep{sarkar2025understanding}. This matters for agentic and
% self-refinement settings, where a model that benefits from good feedback may
% still be a poor teacher or critic. \textsc{ColorRef} makes this measurable
% through directional precision, constraint satisfaction, and paired
% trajectories.

% \paragraph{Behavioral geometry as a complement to representation probing.}
% Prior work asks whether internal representations align with perceptual color
% geometry. \textsc{ColorRef} asks the complementary behavioral question:
% whether outputs and trajectories can be steered toward that geometry. The
% increase in pairwise geometry alignment after feedback suggests that
% interaction reshapes not only individual predictions but also the global
% structure of model behavior, linking grounded semantic evaluation with
% representation-geometry analysis.

% \section{Discussion}
% \label{sec:discussion}

% \paragraph{Interaction as adaptive grounding.}
% The information-budget ablation shows that feedback is not merely an
% additional information channel. At matched total constraint budgets,
% sequential feedback outperforms delivering the same constraints in a single
% turn, suggesting that interaction helps because each teacher message is
% conditioned on the model's previous guess. In this sense,
% \textsc{ColorRef} tests whether language can act as a control signal for
% navigating perceptual space. A correction such as \textit{more blue} is
% useful only relative to the current guess; once the guess changes, the next
% useful correction changes as well.

% \paragraph{Why three turns?}
% The longer-turn sweep supports three feedback rounds as the main protocol.
% Final error improves rapidly from one to three turns and then largely
% plateaus. Longer games still reduce best-so-far error, but the growing gap
% between final and best-so-far error indicates overshooting. Three turns
% therefore provide a practical compute--quality tradeoff, while longer games
% are better suited for studying stopping policies or best-state selection.

% \paragraph{Demonstrations calibrate priors, feedback localizes targets.}
% In-context examples modestly improve one-shot predictions, especially for
% abstract-idiosyncratic descriptions, but add little once oracle feedback is
% available. This suggests that demonstrations calibrate the initial color
% prior, whereas per-example feedback supplies the stronger localization
% signal. The LoRA finetuning result points in the same direction: small-scale
% supervision does not substantially change the conclusions, indicating that
% the main gains come from adaptive feedback rather than format or dataset
% calibration alone.

% \paragraph{Following and generating feedback are distinct capabilities.}
% The LLM-teacher results reveal a sharp asymmetry: the guesser can often
% follow well-formed perceptual corrections, but LLM teachers frequently
% generate fluent feedback that points in the wrong geometric direction. This
% is a controlled form of common-ground misalignment: the teacher produces a
% plausible utterance that is not aligned with the perceptual relation needed
% for update \citep{sarkar2025understanding}. This matters for agentic and
% self-refinement settings, where a model that benefits from good feedback may
% still be a poor teacher or critic. \textsc{ColorRef} makes this failure mode
% measurable through directional precision, constraint satisfaction, and paired
% trajectories.

% \paragraph{Behavioral geometry as a complement to representation probing.}
% Prior work asks whether internal representations align with perceptual color
% geometry. \textsc{ColorRef} asks the complementary behavioral question:
% whether outputs and trajectories can be steered toward that geometry. The
% increase in pairwise geometry alignment after feedback suggests that
% interaction reshapes not only individual predictions but also the global
% structure of model behavior, linking grounded semantic evaluation with
% representation-geometry analysis.


\section{Discussion}
\label{sec:discussion}

\paragraph{Interaction as adaptive grounding.}
The information-budget ablation shows that feedback is not merely an
additional information channel. At matched total constraint budgets,
sequential feedback outperforms delivering the same constraints in a single
turn because later teacher messages are conditioned on the model's previous
guess. In this sense, \textsc{ColorRef} tests whether language can act as a
control signal for navigating perceptual space: a correction such as
\textit{more blue} is useful only relative to the current guess, and changes
as the guess changes.

\paragraph{Why three turns?}
The longer-turn sweep supports three feedback rounds as the main protocol:
final error improves rapidly from one to three turns and then largely
plateaus. Longer games reduce best-so-far error, but the growing final--best
gap indicates overshooting. Three turns therefore provide a practical
compute--quality tradeoff, while longer games are better suited for stopping
policies or best-state selection.

\paragraph{Demonstrations calibrate priors, feedback localizes targets.}
In-context examples modestly improve one-shot predictions, especially for
abstract-idiosyncratic descriptions, but add little once oracle feedback is
available. This suggests that demonstrations calibrate the initial color
prior, whereas per-example feedback supplies the stronger localization
signal. The LoRA result points in the same direction: small-scale supervision
does not substantially change the conclusions, indicating that the main
gains come from adaptive feedback rather than dataset calibration alone.

\paragraph{Following and generating feedback are distinct capabilities.}
The LLM-teacher results reveal a sharp asymmetry: the guesser can often
follow well-formed perceptual corrections, but LLM teachers frequently
generate fluent feedback that points in the wrong geometric direction. This
is a controlled form of common-ground misalignment, where the teacher
produces a plausible utterance that is not aligned with the perceptual
relation needed for update \citep{sarkar2025understanding}. This matters for
agentic and self-refinement settings: a model that benefits from good
feedback may still be a poor teacher or critic. \textsc{ColorRef} makes this
failure mode measurable through directional precision, constraint
satisfaction, and paired trajectories.

\paragraph{Behavioral geometry as a complement to representation probing.}
Prior work asks whether internal representations align with perceptual color
geometry. \textsc{ColorRef} asks whether outputs and trajectories can be
steered toward that geometry. The increase in pairwise geometry alignment
after feedback suggests that interaction reshapes not only individual
predictions but also the global structure of model behavior, linking grounded
semantic evaluation with representation-geometry analysis.


% -------------



\section{Conclusion and Future Work}
\label{sec:conclusion}
We introduced \textsc{ColorRef}, an interactive reference-game framework for
probing language-to-color grounding beyond one-shot prediction. The results
show that controlled feedback improves grounding, bandwidth and sequential
adaptation matter, abstract descriptions benefit strongly from interaction,
and models can follow valid perceptual feedback better than they can generate
it. Future work should extend this approach to multilingual data, human
participants, distributional color targets, and
other continuous perceptual domains.

\section{Limitations}
\label{sec:limitations}

Our study has several limitations. First, all main experiments use English
color descriptions. Extending \textsc{ColorRef} to multilingual data will
require distributional comparisons rather than parallel item-level matching,
since color naming systems and naturally occurring descriptions differ
across languages. Second, the scaled experiments use Qwen3-14B, with
Qwen3-8B as a same-family replication; broader claims require evaluation
across additional model families and multimodal systems. Third, the
LLM-teacher experiments are diagnostic-scale, using \texttt{debug\_4000}, so
they identify a failure mode rather than exhaustively benchmarking teacher
models. Fourth, oracle teachers are controlled probes that provide
geometrically valid feedback by construction; they should not be interpreted
as models of human tutoring. Finally, we evaluate against single
crowd-sourced sRGB targets using Euclidean distance in CIELAB
($\Delta E_{76}$). Abstract descriptions may evoke distributions of
plausible colors across speakers, and future work should incorporate human
judgments and alternative perceptual metrics such as CIEDE2000.


% Bibliography entries for the entire Anthology, followed by custom entries
%\bibliography{anthology,custom}
% Custom bibliography entries only
\bibliography{custom}

\appendix

\section{Prompt Templates}
\label{sec:appendix-prompts}

This section reports the prompt templates used for the guesser and teacher
conditions. All generations use deterministic decoding with temperature
$0$.

\subsection{One-shot and Initial Guess Prompt}

The one-shot baseline and the first turn of feedback games use the same
template.

\begin{quote}
\small
\begin{verbatim}
You are playing a color guessing game.

Color description:
"{raw_name}"

Guess the intended color as a hex code.
Answer only with a valid hex code in the form 
#RRGGBB.
\end{verbatim}
\end{quote}

\subsection{Revision Prompts}

After each teacher message, the guesser receives the original description,
its previous guess, and the latest feedback. The default prompt uses only the
most recent revision step:

\begin{quote}
\small
\begin{verbatim}
You are playing a color guessing game.

Color description:
"{raw_name}"

Your previous guess was: {previous_guess_hex}
Feedback: {feedback}

Revise your guess.
Answer only with a valid hex code in the form 
#RRGGBB.
\end{verbatim}
\end{quote}

We also evaluate a history-aware revision prompt, where the guesser receives
the full interaction history for the current game rather than only the latest
feedback message:

\begin{quote}
\small
\begin{verbatim}
You are playing a color guessing game.

Color description:
"{raw_name}"

Game history:
Turn 0 guess: {guess_0_hex}
Feedback 0: {feedback_0}

Turn 1 guess: {guess_1_hex}
Feedback 1: {feedback_1}

Turn 2 guess: {guess_2_hex}
Feedback 2: {feedback_2}

Revise your guess using the full history above.
Answer only with a valid hex code in the form 
#RRGGBB.
\end{verbatim}
\end{quote}


\subsection{Oracle Feedback Templates}

Oracle teachers compute target--guess differences in LAB lightness
$L^\ast$, LAB red--green $a^\ast$, LAB yellow--blue $b^\ast$, and HSV
saturation. Candidate axes are filtered by minimum thresholds and ranked by
normalized absolute error. The axis oracle emits up to $k$ corrections per
turn, where $k \in \{1,2,3\}$.

Example oracle feedback messages are:

\begin{itemize}
    \item one constraint: \textit{Make it lighter.}
    \item two constraints: \textit{Make it lighter and more blue.}
    \item three constraints: \textit{Make it lighter, more blue, and more muted.}
\end{itemize}

The minimal oracle is the special case $k=1$. The template oracle uses the
same selected axes as the axis oracle, but realizes each direction with a
deterministic paraphrase chosen from a small template set.

\subsection{LLM Teacher Prompts}
\label{ap:teacher_prompts}

The LLM teacher receives the target color and the current guess, and is
asked to provide concise feedback without revealing the target. The prompt
uses a constrained perceptual vocabulary including
\textit{lighter/darker}, \textit{more red/more green},
\textit{more yellow/more blue}, \textit{more saturated/more muted}, and
\textit{warmer/cooler}. In the LAB-aware condition, the teacher additionally
receives target and guess LAB coordinates. In the oracle-assisted condition,
oracle-selected directions are provided and the LLM is asked to paraphrase
them. 

Main LLM teacher used: 
\begin{quote}
\small
\begin{verbatim}
You are the teacher in a color reference game.

Target color description:
"{raw_name}"

True target color: {target_hex}
Current guess: {guess_hex}

Give one short piece of comparative feedback to help 
the guesser move closer to the true target.

Use only these feedback types:
- lighter or darker
- more red or more green
- more yellow or more blue
- more saturated or more muted
- warmer or cooler

Use at most {max_clauses} feedback clause(s) in one 
short sentence.

Do not reveal the target hex code.
Do not mention a new object example.
Do not say the exact answer.
Answer with one short sentence only.
\end{verbatim}
\end{quote}

The Oracle-assisted variant: 
\begin{quote}
\small
\begin{verbatim}
You are the teacher in a color reference game.

Target color description:
"{raw_name}"

True target color: {target_hex}
Current guess: {guess_hex}

The correct perceptual corrections for this step 
are (paraphrase only, do not list as JSON):
{oracle_constraints_json}

Turn these into one natural short sentence of 
comparative feedback for the guesser.

Use only these feedback types:
- lighter or darker
- more red or more green
- more yellow or more blue
- more saturated or more muted
- warmer or cooler
- closer to gray or less gray

Use at most {max_clauses} feedback clause(s).

Do not reveal the target hex code.
Do not mention exact RGB or LAB values.
Do not mention a new object example.
Do not say the exact answer.
Answer with one short sentence only.
\end{verbatim}
\end{quote}

The LAB variant: 
\begin{quote}
\small
\begin{verbatim}
You are the teacher in a color reference game.

Target color description:
"{raw_name}"

True target color: {target_hex}
Target LAB (L, a, b): {target_lab_l}, 
{target_lab_a}, {target_lab_b}

Current guess: {guess_hex}
Current guess LAB (L, a, b): {guess_lab_l}, 
{guess_lab_a}, {guess_lab_b}

Give one short piece of comparative feedback to 
help the guesser move closer to the true target.

Use only these feedback types:
- lighter or darker
- more red or more green
- more yellow or more blue
- more saturated or more muted
- warmer or cooler
- closer to gray or less gray

Use at most {max_clauses} feedback clause(s) in 
one short sentence.

Do not reveal the target hex code.
Do not quote exact LAB numbers in your answer.
Do not mention a new object example.
Do not say the exact answer.
Answer with one short sentence only.


\end{verbatim}
\end{quote}



\subsection{Bandwidth and Wording Conditions}
\label{app:bandwidth_wording_examples}

The controlled oracle conditions vary two factors: feedback
\emph{bandwidth}, the maximum number of constraints emitted in one turn, and
feedback \emph{wording}, the surface form used to express those constraints.
The subscript $c1$, $c2$, or $c3$ denotes the maximum number of constraints
per teacher message.

For example, suppose the target is lighter, bluer, and more muted than the
current guess. The conditions differ as follows:

\begin{center}
\small
\begin{tabular}{@{}ll@{}}
\toprule
Condition & Example feedback \\
\midrule
$\text{minimal}_{c1}$ &
\textit{Make it lighter.} \\
$\text{axis}_{c1}$ &
\textit{Make it lighter.} \\
$\text{axis}_{c2}$ &
\textit{Make it lighter and more blue.} \\
$\text{axis}_{c3}$ &
\textit{Make it lighter, more blue, and more muted.} \\
$\text{template}_{c1}$ &
\textit{Try making it a little lighter.} \\
$\text{template}_{c2}$ &
\textit{It should be lighter and closer to blue.} \\
$\text{template}_{c3}$ &
\textit{Move it lighter, bluer, and a bit more muted.} \\
\bottomrule
\end{tabular}
\end{center}

The \textbf{axis} conditions use canonical direction phrases with fixed
wording, while the \textbf{template} conditions use the same selected
directions but realize them with deterministic natural-language paraphrases.
The \textbf{minimal} condition is a one-constraint oracle used as a simple
baseline; in practice it is equivalent to the one-constraint axis oracle in
our controlled setting, but we keep the name to distinguish it from the
bandwidth sweep.






\section{Additional Results}
\label{sec:appendix-results}

\subsection{Longer-Turn Convergence}
\label{ap:longer-turn}

The main protocol uses up to three feedback rounds. To test whether
performance continues improving with longer games, we run an axis$_{c3}$
turn sweep on \texttt{main\_4000}. Table~\ref{tab:appendix-convergence}
reports final error, best-so-far error, and regression rate.

\begin{table}[h]
\centering
\small
\begin{tabular}{@{}rrrrr@{}}
\toprule
Turns & Final $\Delta E$ & Best $\Delta E$ & Final--best & Reg. rate \\
\midrule
1  & 31.88 & 26.48 & 5.40  & 28.5\% \\
2  & 24.89 & 19.46 & 5.43  & 22.3\% \\
3  & 21.94 & 15.67 & 6.27  & 19.6\% \\
5  & 21.43 & 12.39 & 9.04  & 20.2\% \\
7  & 20.62 & 11.16 & 9.46  & 17.4\% \\
10 & 21.04 & 10.17 & 10.87 & 15.7\% \\
\bottomrule
\end{tabular}
\caption{
Longer-turn convergence sweep for axis$_{c3}$ on \texttt{main\_4000}.
Final error mostly plateaus after three turns, while best-so-far error
continues to improve but is increasingly not retained.
}
\label{tab:appendix-convergence}
\end{table}

The largest gains occur early. Best-so-far error improves by 15.00
$\Delta E$ from turn 0 to 1, 7.02 from turn 1 to 2, and 3.79 from turn 2
to 3. After turn 3, marginal gains fall quickly: 2.03 from turn 3 to 4,
1.25 from turn 4 to 5, and below 1.0 thereafter. This supports using three
turns as a compute--quality tradeoff for final-answer evaluation.

\subsection{In-Context Calibration}
\label{ap:icl-experiments}
We test whether providing examples of color-description/hex pairs improves
the model's color prior. Table~\ref{tab:appendix-icl} reports one-shot and
feedback-game results by number of demonstrations.

\begin{table}[h]
\centering
\small
\begin{tabular}{@{}r@{\hspace{0.2cm}} r@{\hspace{0.2cm}} r@{\hspace{0.2cm}} r@{\hspace{0.2cm}}@{}}
\toprule
$k$ & One-shot $\Delta E$ & Feedback init $\Delta E$ & Feedback final $\Delta E$ \\
\midrule
0 & 41.47 & 41.47 & 21.94 \\
2 & 38.53 & 38.53 & 21.77 \\
4 & 38.71 & 38.71 & 22.00 \\
8 & 39.25 & 39.25 & 22.29 \\
\bottomrule
\end{tabular}
\caption{
In-context calibration on \texttt{main\_1000}. Demonstrations modestly
improve one-shot prediction, but provide little additional benefit once
oracle feedback is available.
}
\label{tab:appendix-icl}
\end{table}

Averaged over $k=2,4,8$, random and regime-matched examples yield similar
one-shot errors, 38.23 and 38.27 respectively, while text-similar examples
are worse at 40.00. This suggests that lexical similarity between color
names does not necessarily imply perceptual similarity. The largest
one-shot ICL gain occurs for abstract-idiosyncratic descriptions, where
mean error drops from 54.21 to 47.92.

\subsection{Information-Budget Ablation}
\label{ap:info-budget}

Table~\ref{tab:appendix-budget-full} reports the full information-budget
ablation on \texttt{main\_12000}. Conditions vary the number of feedback
rounds, maximum constraints per turn, and total constraint budget.

\begin{table}[h]
\centering
\small
\begin{tabular}{@{}lrrrr@{}}
\toprule
Condition & Rounds & $c$/turn & Budget & Final $\Delta E$ \\
\midrule
one-turn$_{c1}$   & 1 & 1 & 1 & 33.81 \\
one-turn$_{c2}$   & 1 & 2 & 2 & 31.98 \\
one-turn$_{c3}$   & 1 & 3 & 3 & 31.88 \\
two-turn$_{c1}$   & 2 & 1 & 2 & 29.71 \\
two-turn$_{c2}$   & 2 & 2 & 4 & 25.11 \\
three-turn$_{c1}$ & 3 & 1 & 3 & 28.13 \\
three-turn$_{c3}$ & 3 & 3 & 9 & 21.94 \\
\bottomrule
\end{tabular}
\caption{
Information-budget ablation. At matched total budgets, iterative feedback
outperforms delivering the same number of constraints in a single turn.
}
\label{tab:appendix-budget-full}
\end{table}

At budget 3, three one-constraint turns reach final error 28.13, while one
three-constraint turn reaches 31.88. At budget 2, two one-constraint turns
reach 29.71, while one two-constraint turn reaches 31.98. Thus, interaction
adds value beyond total feedback quantity.

\subsection{Model-Size Comparison}
\label{ap:model-size}

Table~\ref{tab:appendix-model-size} reports a same-family model-size
comparison on \texttt{main\_4000}.

\begin{table}[h]
\centering
\small
\begin{tabular}{@{}llrr@{}}
\toprule
Model & Condition & Final $\Delta E$ & Rel. imp. \\
\midrule
Qwen3-14B & one-shot & 41.47 & 0.000 \\
Qwen3-14B & axis$_{c3}$ & 21.94 & 0.254 \\
Qwen3-14B & template$_{c1}$ & 29.97 & 0.128 \\
Qwen3-8B & one-shot & 44.18 & 0.000 \\
Qwen3-8B & axis$_{c3}$ & 27.99 & 0.092 \\
Qwen3-8B & template$_{c1}$ & 35.54 & -0.038 \\
\bottomrule
\end{tabular}
\caption{
Model-size comparison. Qwen3-8B shows the same qualitative ordering as
Qwen3-14B, but benefits less from feedback.
}
\label{tab:appendix-model-size}
\end{table}

\section{Dataset Construction Details}
\label{sec:appendix-dataset}

\subsection{Preprocessing Pipeline}
\label{ap:prepro}

We construct \textsc{ColorRef} from a ColorNames-derived corpus of
user-provided color names paired with sRGB hex codes. The pipeline proceeds
in four stages.

\begin{enumerate}
    \item \textbf{Dataset filtering.} We load source color-name entries and
    LLM-derived validity scores. The validity score estimates whether a
    string functions plausibly as a color description. We retain English
    entries with score greater than 0.75 and valid hex codes. The data was analyzed to detect any PII informartion.
    \item \textbf{Canonicalization.} We assign stable example identifiers,
    normalize basic string fields, tokenize descriptions, and compute length
    features.
    \item \textbf{Color conversion.} We convert target hex codes to RGB,
    HSV, and CIELAB coordinates. Main evaluation uses Euclidean distance in
    CIELAB space.
    \item \textbf{Taxonomy construction.} We run a lightweight parser and
    compute component scores for explicitness, prototype anchoring, rarity,
    and abstraction. These scores are used to assign each description to a
    semantic regime.
\end{enumerate}

The final corpus contains 518{,}830 English examples. Each row contains the
raw color description, target hex code, RGB/HSV/CIELAB coordinates, parser
outputs, component scores, regime label, and perceptual bins.

\subsection{Color-name structural parser}
\label{ap:parser}


We assign each raw name a structural parse using a deterministic, lexicon-driven procedure (no model weights). Strings are normalized to lowercase; apostrophes, hyphens, and quotes are mapped to spaces; non-alphabetic characters are removed; runs of whitespace are collapsed. The normalized string is split on spaces into tokens. Head selection scans tokens from right to left: the first token in a fixed head lexicon of basic color terms is taken as the head. Tokens strictly before the head are partitioned into known modifiers (from a fixed modifier lexicon of intensity, temperature, texture, and similar non-head descriptors) versus unknown prefix tokens that are neither heads nor listed modifiers. From this decomposition we assign one of several parse categories: for example, a single head token only, exactly one known modifier with no unknown prefixes , multiple known modifiers and no unknown prefixes , a head with at least one unknown prefix , or no recognized head . We also assign a parse confidence that is high when the prefix is fully explained by the modifier lexicon and decreases when unknown prefixes appear (capped with a floor for compound cases).

\paragraph{Parser v2 (used downstream)}.
For experiments that require broader coverage, we use an extended parser that applies token-level rules before lexicon lookup: $-i$sh normalization (e.g., mapping greenish to green, with a small exception table), optional plural stripping when the singular stem is a known head or prototype, and—if no basic color head is found—a second right-to-left pass that allows prototype heads (object-, material-, or nature-associated color words from a separate lexicon). That yields additional categories such as prototype-only heads, modifiers with a prototype head, and compounds around a prototype head. Regardless of whether the final head is basic or prototypical, v2 also records all prototype lexicon hits anywhere in the name as auxiliary candidates for later taxonomy features. All lexicons are hand-curated YAML lists checked into the repository.


\subsection{Validity Filtering}
\label{ap:val-filtering}

The validity filter removes entries judged unlikely to function as color
descriptions. We used GPT 4o as a judge with the following prompt (extract): 

\begin{quote}
\small
\begin{verbatim}
You are evaluating whether text entries 
are plausible color descriptions. 
You will receive a list of descriptions 
and must score each one from 0 to 1, where:

- **1.0** = Definitely a plausible color 
description
- **0.5** = Uncertain / borderline case
- **0.0** = Not a meaningful color 
description

## What to ACCEPT (score more than 0.7)

**Literal color descriptions**: "dark red", 
"light blue", "pale yellow"
**Figurative/associative descriptions**: 
"sunset", "ocean", "grandmother's kitchen", 
"tired Monday morning"
**Subjective/abstract descriptions**: 
"melancholy", "peaceful", "nostalgic", 
"hopeful"

## What to REJECT (score less 0.3)

**Spam/promotional**: "buy now", 
"visit my site", "subscribe", "click here"
**URLs/emails**: "www.example.com", 
"user@email.com", any web links
**Nonsensical text**: "asdfgh",
"qwerty", "zzzzz", random characters
**Non-descriptive statements**: "im a genius", 
"hello world", "test 123"
**Personal messages**: "love you mom",
"happy birthday", "call me"

## Borderline Cases (score 0.4-0.6)

**Ambiguous single words**: Could be a 
name or a description? Use context.
**Very abstract concepts**: "philosophy", 
"mathematics" (unless clearly sensory)
**Proper names**: "Jennifer", "Tokyo" 
(unless commonly used for colors)
**Mixed content**: Partially meaningful 
but with errors/spam elements


\end{verbatim}
\end{quote}



The threshold is fixed at score $>0.75$. This filter is used
to remove obvious noise while preserving creative and abstract color names.
The score is not used as a target variable in the experiments.








\section{AI Assistance Disclosure}
\label{app:ai_assistance}

We used AI coding assistance, specifically Claude Code, to automate routine
experiment-management tasks, including launching ablation runs, organizing
outputs, and checking generated result files. All experimental designs,
metrics, analyses, and paper claims were specified and reviewed by the
authors.

\end{document}
